\section{Geo-projected Ground Motion (GGM)}

\subsection{Product Description}
The Geo-projected Ground Motion (\acs{ggm}) product is particularly valuable for ground deformation applications. Its purpose is to derive three-dimensional displacement measurements by combining information from two (or more, when available) acquisition geometries. Unlike conventional SAR-based measurements expressed in the sensor reference frame (azimuth and slant range), the \acs{ggm} expresses motion in a geographic reference system, enabling more intuitive and physically meaningful analysis.

This geo-projected 3D displacement information is especially useful in a wide range of applications, such as glacier motion tracking, earthquake-induced surface deformation, subsidence monitoring, and other geophysical phenomena.

The definition of the \acs{ggm} product involves several critical aspects that must be carefully assessed to determine whether its implementation is justified. One feasible approach is to derive the \acs{ggm} by combining \acs{gpro} \& \acs{gpao} obtained from complementary acquisition geometries. In particular:

\begin{itemize}
    \item \acs{gpro} and \acs{gpao} extracted from phase information. In this case, the product \acs{gpro} coincides with \acs{gunw}. 
    \item \acs{gpro} and \acs{gpao} extracted from amplitude information. 
\end{itemize}
Both of the previous products provide displacement measurements, relative to a reference position, between two acquisition dates in the azimuth/slant range directions by using phase information. By using data from both ascending and descending geometries over the same geographic area, it is possible to compute the three components (north, east, vertical) of the displacement vector in a geographic coordinate frame.

The displacement captured by \acs{ggm} is the cumulative displacement measurements between two relatively close dates. \autoref{tab:ggm_product_characteristics} summarizes the different input configurations and the corresponding displacement components derivable from them.

\begin{longtblr}[
  label = tab:ggm_product_characteristics,
  caption = {Types of Measurements for Geo-projected Ground Motion},
]{
  width = \linewidth,
  colspec = {Q[l,25]Q[l,20]Q[l,20]Q[l,35]},
  hlines,
  hline{1-Z} = {-}{0.05em},
  rowhead=2,
}
\textbf{Input Geometries} & \textbf{GGM Type} & \textbf{Inputs} & \textbf{Measurement Components} \\
\SetCell[r=2]{l} Double geometry (ascending \& descending) & \acs{nev}-\acs{pt}\_GGM & \acs{gpro} \& \acs{gpao} by using amplitude & Along north, east, and vertical directions \\
 & \acs{nev}-IN\_GGM & \acs{gpro} \& \acs{gpao} by using phase & Along north, east, and vertical directions \\
\end{longtblr}

\subsubsection{Input and Auxiliary Data}
The input and auxiliary data  considered for product generation are listed in \autoref{tab:in_aux_ggm_products} (refer to \autoref{tab:ggm_product_characteristics} for \acs{ggm} product type).

\begin{longtblr}[
	label = tab:in_aux_ggm_products,
	entry = none,
	caption = {\acs{ggm} input and auxiliary data},
	]{
		width = \linewidth,
		colspec = {Q[l,20]Q[l,60]Q[c,20]},
		hlines,
		hline{1-Z} = {-}{0.05em},
	}
	\textbf{Name}   & \textbf{Description}   & \textbf{Mandatory (M) / Optional (O)} \\
	\acs{gpro}\_asc & Ground-projected Precise Range Offsets derived from ascending acquisitions by exploring phase or amplitude information & M \\
    \acs{gpro}\_desc & Ground-projected Precise Range Offsets derived from descending acquisitions by exploring phase or amplitude information & M \\
\end{longtblr}

The \acs{ggm} product is generated from input data derived using the same methodology. Therefore, it is not possible to combine pixel offset products obtained through amplitude cross-correlation with those derived from phase information, due to the significantly different sensitivity and accuracy of the input measurements.  

\subsubsection{Main Variables}
\autoref{tab:var_ggm_products} lists the main variables associated with the \acs{ggm} products.

\begin{longtblr}[
	label = tab:var_ggm_products,
	entry = none,
	caption = {\acs{ggm} main variables},
	]{
		width = \linewidth,
		colspec = {Q[l,15]Q[l,52]Q[l,20]Q[l,13]},
		hlines,
		hline{1-Z} = {-}{0.05em},
	}
	\textbf{Name} & \textbf{Description} & \textbf{Dimensions} & \textbf{Data Type} \\
    displacement
    
    \_north & Ground displacement along the north (Y) direction, in millimeters or centimeters. Present in \acs{nev} \acs{ggm} products & x, y & float32 \\
    displacement
    
    \_east & Ground displacement along the east (X) direction, in millimeters or centimeters. Present in \acs{nev} \acs{ggm} products & x, y & float32 \\
    displacement
    
    \_vertical & Ground displacement in the vertical (Z) direction, in millimeters or centimeters. Present in \acs{nev} \acs{ggm} products & x, y & float32 \\
\end{longtblr}

\autoref{tab:var_ggm_coordinates} describes the dimensions/coordinates associated with the \acs{ggm} main variables.

\begin{longtblr}[
    label = tab:var_ggm_coordinates,
    entry = none,
    caption = {\acs{ggm} main variables coordinates/dimensions},
    ]{
        width = \linewidth,
        colspec = {Q[l,20]Q[l,62]Q[l,18]},
        hlines,
        hline{1-Z} = {-}{0.05em},
    }
    \textbf{Name} & \textbf{Description} & \textbf{Data Type} \\
    x & UTM/UPS x (East) coordinates [m] & float64 \\
    y & UTM/UPS y (North) coordinates [m] & float64 \\
\end{longtblr}

\subsubsection{Main Metadata}
The following tables list the main metadata and per-pixel metadata associated with the \acs{ggm} products.

\begin{longtblr}[
	label = tab:meta_ggm_products,
	entry = none,
	caption = {\acs{ggm} main metadata and per-pixel metadata},
	]{
		width = \linewidth,
		colspec = {Q[l,17]Q[l,38]Q[l,20]Q[l,13]Q[c,12]},
		hlines,
		hline{1-Z} = {-}{0.05em},
	}
	\textbf{Name} & \textbf{Description} & \textbf{Dimensions} & \textbf{Data Type} & \textbf{Per-pixel Metadata}\\
    inputs\_IDs & List of the input products used fore generating the \acs{ggm} product & - & List of strings & \\
    date\_start & First or reference date of the interval to which the measurements refer & - & string & \\
    date\_end & Last date of the interval to which the measurements refer & - & string & \\
    reference\_point & ID of the reference point/cell of the measurements & - & string & \\
    quality\_mask & Mask indicating valid/incoherent areas & x, y & byte & \ding{51} \\
    uncertainty
    
    \_north & Uncertainty of the north displacement component. Present in \acs{nev} \acs{ggm} products. & x, y & float32 & \ding{51} \\
    uncertainty
    
    \_east & Uncertainty of the east displacement component. Present in \acs{nev} \acs{ggm} products. & x, y & float32 & \ding{51} \\
    uncertainty
    
    \_vertical & Uncertainty of the vertical displacement component. Present in \acs{nev} \acs{ggm} products. & x, y & float32 & \ding{51} \\
    coherence & Interferometric or amplitude coherence & x, y & float32 & \ding{51} \\
    x\_sample & Pixel spacing in x direction & - & float32 & \\
    y\_sample & Pixel spacing in y direction & - & float32 & \\
\end{longtblr}

The table below describes the dimensions/coordinates associated with the \acs{ggm} main metadata.

\begin{longtblr}[
    label = tab:var_ggm_metadata_coordinates,
    entry = none,
    caption = {\acs{ggm} main metadata coordinates/dimensions},
    ]{
        width = \linewidth,
        colspec = {Q[l,20]Q[l,62]Q[l,18]},
        hlines,
        hline{1-Z} = {-}{0.05em},
    }
    \textbf{Name} & \textbf{Description} & \textbf{Data Type} \\
    x & UTM/UPS x (East) coordinates [m] & float64 \\
    y & UTM/UPS y (North) coordinates [m] & float64 \\
\end{longtblr}

\subsubsection{Reference System Tiling and Slicing}
\autoref{tab:ggm_ref_system_slicing} reports the reference system, the slicing, and grid for \acs{ggm} products.

\begin{longtblr}[
    label=tab:ggm_ref_system_slicing,
    entry = none,
    caption = {Reference system and Slicing for \acs{ggm}},
    ]{
        width=\linewidth,
        colspec = {Q[l,15]Q[l,35]Q[l,35]Q[l,15]},
        %vline{2-5} = {-}{},
        %vline{5-5} = {-}{},
        hlines,
        hline{1-Z} = {-}{0.05em},
        rowhead=1
    }
    \textbf{Product} & \textbf{Reference System} & \textbf{Slicing} & \textbf{Spacing}\\
    \textbf{\acs{ggm}} & Cartographic: UTM/UPS & US-\acs{mgrs} 100 km x 100 km & 50 m \\
\end{longtblr}

\subsection{Product Requirements}
The main problem related to such a product is that more than one acquisition has to be combined and this implies that the optimal (or at least reasonable) couplings among a range of possibilities has to be found. Moreover, in general, it is not always possible to establish what is the “reasonable” coupling without taking into account what the specific use case is (i.e. the best coupling can change when we change the use case the \acs{ggm} has been thought for). One additional negative consequence directly related to the data coupling requirement is that \acs{ggm} production flow is much more complex than the ones designed for other \acs{ard} products. 

Although these considerations lead us to believe that these products are too complex for a routine generation process; we could think of a different approach in which we devolve to the end-user the burden of the data coupling (since is the user who knows the application the GGM will be used for), and we could release as open-source the code necessary for the \acs{ggm} production starting from the \acs{ard}s selected by the user.

\subsubsection{Accuracy}

Before assessing the accuracy performances for \acs{ggm} products, a series of consideration must be done:

The accuracy of Geoprojected Ground Motion (\acs{ggm}) products is primarily dictated by the directional sensitivity and noise characteristics of the input measurements:
\begin{itemize}
    \item \acs{los} (Line of Sight) phase measurements provide high precision and are predominantly sensitive to vertical and east-west deformation components, but exhibit limited sensitivity to north-south motion due to the near-polar orbit geometry. 
    \item In contrast, \acs{lof} (Line of Flight) offsets offer stronger sensitivity to the north-south component but are significantly noisier, typically by an order of magnitude or two, than range offsets, owing to lower coherence and reduced measurement precision.
\end{itemize}
In light of these constraints, the approach evaluated for \acs{ggm} products generation is a full 3D retrieval method that combines azimuth and range offsets from both ascending and descending acquisitions, yielding four independent measurements that jointly constrain the vertical, east, and north components of surface displacement. 

Accuracy assessment is conducted with consideration of the measurement geometries, noise levels, and the propagation of uncertainties through the inversion process. 

The following tables summarize the resulting estimated output accuracies for both \acs{s1} and \acs{rosel} missions. 

The following assumptions (see \autoref{tab:pro_pao_accuracy_phase_s1}) have been used to derive the performances of the \acs{s1} mission:  
\begin{itemize}
\item Range offset accuracy $\sigma_{\Delta^{ph}_{rg}} \simeq 0.005$ m 
\item Azimuth offset accuracy $\sigma_{\Delta^{amp}_{az}} \simeq 0.04$ m. 
\end{itemize}

\begin{longtblr}[
  label = tab:ggm_accuracy_s1,
  entry = none,
  caption = {Accuracy Requirements for \acs{ggm} products (\acs{s1} mission)},
]{
  width = \linewidth,
  colspec = {Q[l,35] Q[l,55] Q[c,20]},
  hlines,
  hline{1-Z} = {-}{0.05em},
  rowhead = 1,
}

\textbf{Metric} & \textbf{Description} & \textbf{Requirement} \\

$\sigma_n$ & North-South component accuracy & $\leq 0.03$ m \\

$\sigma_e$ & East-West component accuracy & $\leq 0.006$ m \\

$\sigma_u$ & Vertical component accuracy & $\leq 0.006$ m \\

\end{longtblr}

The following assumptions (see \autoref{tab:pro_pao_accuracy_phase_s1}) have been used to derive the performances for the \acs{rosel} mission:  
\begin{itemize}
\item Range offset accuracy $\sigma_{\Delta^{ph}_{rg}} \simeq 0.02$ m 
\item Azimuth offset accuracy $\sigma_{\Delta^{amp}_{az}} \simeq 0.12$ m. 
\end{itemize}

\begin{longtblr}[
  label = tab:ggm_accuracy_rosel,
  entry = none,
  caption = {Accuracy Requirements for \acs{ggm} products (\acs{rosel} mission)},
]{
  width = \linewidth,
  colspec = {Q[l,35] Q[l,55] Q[c,20]},
  hlines,
  hline{1-Z} = {-}{0.05em},
  rowhead = 1,
}

\textbf{Metric} & \textbf{Description} & \textbf{Requirement} \\

$\sigma_n$ & North-South component accuracy & $\leq 0.09$ m \\

$\sigma_e$ & East-West component accuracy & $\leq 0.025$ m \\

$\sigma_u$ & Vertical component accuracy & $\leq 0.025$ m \\

\end{longtblr}

\subsubsection{Timeliness}
This is a complex product derived from the combination of multiple \acs{ard} products. The input could derive from interferometric unwrapped phases \acs{gunw}, or precise azimuth and range offsets \acs{gpro} and \acs{gpao}. The combination of the different information does not significantly impact timeliness.

\subsection{Usability, Integrability and Interoperability}
The characteristics of the \acs{ggm} products in terms of usability, integrability and interoperability are in the following table.

\begin{longtblr}[
  caption = {Usability, Integrability, and Interoperability of the \acs{ggm} Product},
  label = {tab:ggm_usability_integrability},
]{
  colspec = {Q[2cm] Q[12cm]},
  rowhead = 1,
  width = \linewidth
}
\hline
\textbf{Aspect} & \textbf{Description} \\
\hline
Usability &
\begin{itemize}
  \item 3D Displacement Estimation (in case of multi-geometry derivation): Provides north-east-vertical(\acs{nev}) ground motion, which is easier to interpret than SAR-native (azimuth/slant-range or \acs{los}) measurements.
  \item Application-Ready: Useful for deformation studies related to earthquakes, landslides, subsidence, glacier flow, and infrastructure monitoring.
  \item User-Friendly Geometry: Expressed in geographic coordinates, facilitating integration into \acs{gis} and downstream geophysical analysis.
  \item Take into account that the product accuracy may depend on assumptions (e.g., negligible north-south motion in \acs{gunw}-only cases), and input data quality (e.g., coherence, ionospheric corrections).
\end{itemize}
\\
Integrability &
\begin{itemize}
  \item It Inherits integrability features from \acs{gpro}, \acs{gpao}, and \acs{gunw}, allowing seamless use in existing deformation workflows.
  \item Integration with \acs{gis} and Modeling Tools: Usable as input in hazard models, crustal and glacier deformation simulation or ground deformation assessments, in general.
  \item Suitable for Multi-Source Analysis: Can be fused with other data for validation or joint deformation analysis.
\end{itemize}
\\
Interoperability &
\begin{itemize}
  \item It Uses standard geographic grids and reference systems (e.g., WGS84, \acs{utm}), supporting integration across \acs{eo} systems.
  \item Cross-Mission Compatibility: It can be built from different SAR missions data.
  \item Format Compatibility: Delivered in common geospatial formats, compatible with \acs{eo} tools.
  \item Supports Interdisciplinary Use: Usable by geologists, civil engineers, glaciologists, and emergency response planners.
\end{itemize}
\\
\hline
\end{longtblr}
